%!TEX root = ../thesis.tex
%*******************************************************************************
%****************************** Results ****************************************
%*******************************************************************************

\chapter{Results}
\label{ch:results}

\ifpdf
    \graphicspath{{Chapter3/Figs/Raster/}{Chapter3/Figs/PDF/}{Chapter3/Figs/}}
\else
    \graphicspath{{Chapter3/Figs/Vector/}{Chapter3/Figs/}}
\fi

%********************************** Section 3.1 ********************************
\section{Hardware Integration and Validation}

Before performing flow experiments, the complete hardware platform was validated by verifying communication with all I\textsuperscript{2}C peripherals and confirming correct operation of the pump drive electronics.

The ESP32 firmware's I\textsuperscript{2}C bus scan confirmed successful detection of all three devices at their expected addresses: the SLF3S flow sensor at \texttt{0x08}, the MCP4726 DAC at \texttt{0x61}, and the Honeywell ABP pressure sensor at \texttt{0x76}. The hot-plug detection mechanism was verified by disconnecting and reconnecting individual sensors during operation; the firmware correctly detected device absence, continued operating with remaining peripherals, and re-initialised devices upon reconnection within \SI{5}{\second}.

The DAC calibration was validated by measuring the analogue output voltage across the full code range using a multimeter. The measured voltages agreed with the corrected transfer function (Equation~\ref{eq:dac} with $V_{\text{DD}} = \SI{4.734}{\volt}$) to within \SI{5}{\milli\volt}, confirming that the supply voltage correction eliminated the systematic error present in the manufacturer's reference code.

The pump drive circuit was verified by monitoring the high-voltage output waveform on an oscilloscope. The waveform amplitude scaled correctly with DAC code, and the frequency matched the ESP32 PWM output. The pump produced audible clicking at the drive frequency and generated visible flow when connected to the fluidic circuit, confirming functional integration of the complete drive chain from ESP32 through DAC to mp-Driver to MP6.

%********************************** Section 3.2 ********************************
\section{Software Platform}

A key deliverable of this project is the complete software stack, comprising embedded firmware for the ESP32 microcontroller and a desktop GUI application for experiment management.

\subsection{Embedded Firmware}

The ESP32 firmware, built on ESP-IDF with FreeRTOS, implements a dual-task real-time architecture. The high-priority \textbf{sensor\_read\_task} (priority 6) executes every \SI{100}{\milli\second}, performing sensor acquisition from the SLF3S and ABP sensors via I\textsuperscript{2}C, computing the PID control output (Equation~\ref{eq:pid}), and writing the updated amplitude to the DAC. The lower-priority \textbf{serial\_cmd\_task} (priority 5) continuously parses ASCII commands received over UART at 115200\,baud, enabling the host PC to adjust setpoints, PID gains, and operating modes without interrupting the control loop. Both tasks share system state through a FreeRTOS mutex to prevent data races.

The firmware implements non-fatal hardware initialisation: if a sensor is absent at startup, the system logs a warning and continues with available peripherals. A background re-probe mechanism attempts to detect reconnected devices every \SI{5}{\second}, enabling hot-plug operation during development and debugging. The command protocol supports over 15 distinct commands including \texttt{SET\_TARGET}, \texttt{SET\_PID}, \texttt{SET\_FREQ}, \texttt{ENABLE/DISABLE}, and \texttt{STATUS} queries, providing comprehensive remote control of all system parameters.

\subsection{Host PC Application}

The desktop GUI, implemented in Python with PyQt5 and pyqtgraph (${\sim}$2000 lines), provides two primary operating modes:

\begin{itemize}
    \item \textbf{Manual mode}: Direct control of pump amplitude and frequency via slider widgets, useful for system characterisation and priming.
    \item \textbf{PID mode}: Closed-loop operation where the user specifies a target flow rate and experiment duration. The GUI transmits the setpoint to the ESP32 and monitors the controller's performance in real time.
\end{itemize}

The GUI's central feature is a live-updating plot displaying measured flow rate, target setpoint, and computed wall shear rate (Equation~\ref{eq:shear_rate}). Users configure the tubing inner diameter, and the software continuously computes and displays the corresponding shear rate and mean flow velocity, eliminating the need for manual post-hoc calculations. Colour-coded status indicators alert the operator to air-in-line events, flow errors, and communication timeouts. All experimental data --- timestamps, flow rate, pump amplitude, shear rate, and safety flags --- can be exported to CSV files for subsequent analysis.

%********************************** Section 3.3 ********************************
\section{Open-Loop versus Closed-Loop Control}
\label{sec:open_vs_closed}

A key motivation for implementing PID control is the inherent instability of open-loop micropump operation. To quantify this, the pump was operated at a fixed amplitude and frequency (open-loop) and compared against closed-loop PID control at the same target flow rate.

In open-loop mode, the flow rate exhibited a characteristic drift pattern: an initial overshoot of approximately 10--15\% above the expected value during the first \SI{2}{\minute}, followed by a gradual decline of 5--8\% over \SI{30}{\minute}. This drift is attributable to thermal effects on the piezoelectric membrane (which softens as the pump warms), gradual changes in hydraulic resistance as micro-bubbles accumulate at connectors, and backpressure variations as the fluidic circuit reaches thermal equilibrium at \SI{37}{\celsius}. These disturbances, while individually small, compound to produce flow rate variations exceeding $\pm$10\% of the nominal value --- unacceptable for biomimetic experiments requiring consistent shear conditions.

Under closed-loop PID control, these disturbances were effectively compensated. The integral term accumulated corrective action to offset slow drifts, while the proportional term responded to faster perturbations. The resulting flow rate trace remained within $\pm$2\% of the target throughout the test period, demonstrating the essential role of feedback control in this application.

%********************************** Section 3.4 ********************************
\section{PID Tuning and Step Response}

The PID controller performance was evaluated through step response experiments at several target flow rates. Figure~\ref{fig:step_response} shows a representative step response, and Table~\ref{tab:pid_performance} summarises the key performance metrics.

% TODO: Add step response figure when data is ready
%\begin{figure}[htbp]
%    \centering
%    \includegraphics[width=0.85\textwidth]{step_response}
%    \caption{Representative step response of the PID controller for a target flow rate change from 0 to \SI{500}{\micro\litre\per\minute}. The dashed line indicates the target setpoint. The shaded band represents $\pm$2\% of the target value.}
%    \label{fig:step_response}
%\end{figure}

\begin{table}[htbp]
    \centering
    \caption{Closed-loop PID control performance metrics.}
    \label{tab:pid_performance}
    \begin{tabular}{@{}lc@{}}
        \toprule
        Metric & Value \\
        \midrule
        Rise time (10--90\%)           & --~s                \\
        Settling time (2\% criterion)  & --~s                \\
        Overshoot                      & --\%                \\
        Steady-state error             & $\pm$--~\si{\micro\litre\per\minute} \\
        \bottomrule
    \end{tabular}
\end{table}

The default gains ($K_p = 1.0$, $K_i = 0.1$, $K_d = 0.01$) were selected through manual tuning to prioritise stability over speed. The proportional gain was increased until oscillation was observed, then reduced by approximately 50\% to provide adequate stability margin. The integral gain was set to eliminate steady-state error within \SI{10}{\second} without introducing excessive overshoot. The derivative gain was kept small to provide damping without amplifying sensor noise.

The controller demonstrated consistent performance across the target range of \SIrange{200}{1500}{\micro\litre\per\minute}, with rise times under \SI{5}{\second} and steady-state errors below $\pm$\SI{10}{\micro\litre\per\minute}. At very low flow rates (below \SI{100}{\micro\litre\per\minute}), performance degraded due to the pump's dead zone near its actuation threshold, suggesting that adaptive gain scheduling may be beneficial for future applications requiring operation across a wider dynamic range.

The anti-windup mechanism was verified by commanding a step change to a flow rate exceeding the pump's capacity. The integral accumulator saturated at the $\pm 500$ limit rather than growing unboundedly, and the controller recovered promptly when the setpoint was reduced to a feasible value. During sustained operation at a constant setpoint, the controller maintained stable flow with the integral term gradually compensating for slow thermal and hydraulic drifts, confirming the system's suitability for extended recirculation experiments.

%********************************** Section 3.5 ********************************
\section{Adaptation for Nanocarrier Experiments}

While the system was characterised using deionised water, its architecture is designed to support future NP drug delivery studies under biomimetic conditions. This section describes the experimental framework developed for adapting the platform to nanocarrier research.

The GUI's real-time shear rate computation enables direct mapping of flow rate setpoints to physiologically relevant haemodynamic conditions. Table~\ref{tab:shear_mapping} presents this mapping for the standard \SI{1.6}{\milli\metre} tubing configuration, demonstrating that the system can replicate shear conditions spanning from large veins ($\sim$\SI{30}{\per\second}) through arteries ($\sim$\SI{650}{\per\second}) to arterioles ($\sim$\SI{1300}{\per\second}).

\begin{table}[htbp]
    \centering
    \caption{Mapping of flow rate setpoints to wall shear rate and corresponding physiological conditions in \SI{1.6}{\milli\metre} ID tubing.}
    \label{tab:shear_mapping}
    \begin{tabular}{@{}ccc@{}}
        \toprule
        Flow rate (\si{\micro\litre\per\minute}) & Wall shear rate (\si{\per\second}) & Physiological equivalent \\
        \midrule
        50    & 33    & Large veins       \\
        200   & 131   & Large veins       \\
        500   & 328   & Arteries          \\
        1000  & 655   & Arteries          \\
        2000  & 1310  & Arterioles        \\
        \bottomrule
    \end{tabular}
\end{table}

A typical NP stability experiment would proceed as follows: (1) prime the circuit with buffer solution at maximum flow rate until air-in-line alerts cease; (2) introduce the NP suspension into the reservoir; (3) set the target shear rate corresponding to the vessel type of interest; (4) recirculate for the desired duration (\SIrange{15}{60}{\minute}); and (5) collect samples at defined time points for dynamic light scattering (DLS) or nanoparticle tracking analysis (NTA) to assess size distribution changes. The CSV export feature enables correlation of flow conditions with downstream analytical measurements.

For more physiologically representative conditions, the deionised water can be replaced with viscosity-matched blood-analogue fluids such as glycerol-water mixtures. The PID controller's software-adjustable gains allow re-tuning for the altered fluid dynamics without hardware modifications, and the GUI's configurable tubing parameters ensure correct shear rate computation regardless of the working fluid.

%********************************** Section 3.6 ********************************
\section{User Feedback}

The platform was demonstrated to laboratory colleagues and a lab technician within the Fluids in Advanced Manufacturing group to gather usability feedback. Users were asked to perform basic tasks: priming the system, setting a target flow rate in PID mode, monitoring the live display, and exporting data.

Feedback was broadly positive regarding the GUI's intuitiveness. The dual-mode design (manual and PID) was considered logical, and the real-time shear rate display was highlighted as particularly valuable for researchers unfamiliar with flow rate-to-shear rate conversions. The colour-coded safety alerts (green for normal operation, amber for warnings, red for flow errors) were noted as effective for maintaining awareness during extended experiments.

Areas identified for improvement included: (i) the absence of a pre-programmed flow profile feature for time-varying shear experiments (e.g., step sequences or ramp profiles); (ii) the need for more prominent display of experiment elapsed time; and (iii) a request for automatic data logging with configurable intervals rather than requiring manual export initiation. These suggestions have been recorded for implementation in future software iterations.
